% Source-derived chapter 04: Pairings
% Original source: universal-p-adic-gross-zagier-formula
\chapter{Pairings}
\label{universal-p-adic-gross-zagier-formula-chapter-04}
\source{universal-p-adic-gross-zagier-formula}

\begin{remark}[Source section]
\label{universal-p-adic-gross-zagier-formula-chapter-04-source}
\source{universal-p-adic-gross-zagier-formula}
\source{universal-p-adic-gross-zagier-formula}
This chapter reproduces the corresponding section of the retrieved
LaTeX source, including its definitions, statements, examples, and
proofs. It has not yet been translated into Lean.
\notready
\end{remark}

% The source section title was: Pairings
\subsection{Global dualities}  \label{sec: duality} We construct Hecke- and/or Galois-equivariant duality pairings on the sheaves constructed in the previous section. The results of this somewhat technical subsection are summarised in Propositions \ref{duality}, \ref{abovco}. 
\source{universal-p-adic-gross-zagier-formula}

\subsubsection{Pairings, symmetries and involutions}\lb{skew def}
If $\epsilon \in \{\pm 1\}$,  $R$ is a ring or sheaf of rings, and $M$, $J$ are $R$-modules, an $R$-bilinear pairing $(\ ,  \ )\colon M\ot M\to J$ is said to be \emph{$\epsilon$-symmetric}  if it satisfies $(m, m')=\epsilon\cd (m', m)$. If $R$ is equipped with an involution $\iota$,  denote by  $(\cd)^{\iota}$ both the functor $\ot_{R, \iota}R$ and the maps $m\mapsto m\ot1$;  an $(R, \iota)$-sesquilinear pairing $(\ ,  \ )\colon M\ot M^{\iota}\to J$ is said to be \emph{$\epsilon$-hermitian} if it satisfies $  (m, n)= \epsilon \cd \iota ((n^{\iota}, m^{\iota})^{\iota})$.   We will also use the prefix `skew-' (respectively no prefix) if $\epsilon=-1$ (respectively $+1$). 



\subsubsection{Involutions}
We denote by  the same name $\iota$  the involutions on $\cH_{\G_{*}}^{\rm sph}$, $\cH_{\G_{*}}^{\rm sph, \ord}$ $\Lambda_{\G_{*}, U^{p}_{*}} $,  $\cE_{\G_{*}}^{\ord}$ deduced from those of  \S~\ref{ssec invo}. 
If $M$ is a module over any of the above rings (or sheaf of modules over any of the above spaces), we let 
$M^{\iota}= \iota^{*} M$.  

\begin{lemm}\label{iota and dual} Let $W$ be an irreducible   algebraic representation of $\G_{*}$ over $L$.
\source{universal-p-adic-gross-zagier-formula}
\begin{enumerate}
\item We have $\sigma_{W^{\vee}}(t)=\sigma_{W}(t^{\iota})$ for all $t\in T_{\G_{*}}$.
\item
If $\pi^{\ord}$ is the ordinary part of  an automorphic representation of $\G_{*}(\A^{\infty})$ over $L$ of weight $W$, unramified of level $U_{*}^{S}$ outside of a finite set of primes $S$, then there is an isomorphism of $\cH_{\G_{*}, U^{S}}^{{\rm sph}, \ord}$-modules $\pi^{\vee,U_{*}^{S}, \ord}\cong (\pi^{U_{*}^{S},\ord})^{\iota}$.
  \item There is an identification 
$$(\cE_{G_{*}}^{\cl, W})^{\iota}=\cE_{\G_{*}}^{\cl, W^{\vee}}$$
such that $\pi^{\ord}_{\iota(x)}=(\pi_{x})^{\vee, \ord}$. 
\end{enumerate}
\end{lemm}
\begin{proof} All results can be reduced to the case $\G_{*}=\H$, which is trivial, or $\G_{*}=\G$, that we address. Part 1 follows from the explicit description of $\sigma_{W}$ in \eqref{alg char} and $W_{\G, (w; (w_{\sigma}))}^{\vee}\cong W_{\G,(-w; (w_{\sigma}))}$ (see \eqref{inv pair W} below for an explicit duality). 

For part 2, we use $\pi^{\vee}=\pi\otimes \omega^{-1}$ where $\omega$ is the central character of $\pi$, and verify the statement separately for the spherical Hecke algebra and for the operators $\Up_{t}$. For the former, it is well known that the spherical Hecke algebra is generated by operators $T(z)$ and $T(\smalltwomat x{}{}1)=T(\smalltwomat 1{}{}x)$ for $z,x\in F_{S}^{\times}$; denoting by $\lambda_{\pi^{?}}(\cdot) $ the eigenvalue of $T(\cdot) $ on $(\pi^{?})^{U^{S}}$, we then have 
$\lambda_{\pi}(z^{\iota})=\lambda_{\pi}(z^{-1})=\omega(z)^{-1}=\lambda_{\pi^{\vee}}(z)$, and $$\lambda_{\pi}((\smalltwomat x{}{}1^{\iota})= \lambda_{\pi}(x^{-1}(\smalltwomat 1{}{}x)=\omega(x)^{-1}\lambda_{\pi}(\smalltwomat x{}{}1)=\lambda_{\pi^{\vee }}(\smalltwomat x{}{}1)$$ as desired. 

For the operators $\Up_{t}$, we verify that if $\pi$ is ordinary at $v$ with unit character $\alpha_{v}^{\circ}=\alpha_{v}\sigma_{W}^{-1}$ (as a character of $T_{\G, v}^{+}$),  then $\pi^{\vee}$ is ordinary at $v$ with unit character 
$$\alpha_{v}^{\circ,\iota}\colon t\mapsto \alpha_{v}^{\circ} (t^{\iota}).$$
This follows from observing 
\beqq
 \pi_{v}^{\vee} & \cong {\rm Ind} ( \alpha_{v}^{ \iota}\cdot (|\  |_{v},|\  |_{v}^{-1})),\\
\alpha_{v}^{\circ}(t^{\iota}) & =\alpha_{v}^{\circ}(t)  \alpha_{v}^{\circ}(\nu(t))^{-1}\in \OO_{F}^{\ts}.
\eeqq




Finally, part 3 follows from parts 1 and 2. 
\end{proof}








\subsubsection{Homological and cohomological  dualities} 
We shall define various pairings $\la\ , \  \ra_{?}$ in the (ordinary, completed) homology of Shimura varieties, starting from the Poincar\'e duality pairings. Then we will use them to construct corresponding pairings $(\,  \ )_{?}$ on spaces of representations, as follows.


\begin{enonce}{Construction}\label{constr pairing} 
\source{universal-p-adic-gross-zagier-formula}
Let $A$ be a ring, $G$ a group,  and let $M_{1}, M_{2}, V_{1}, V_{2}, A(d)$ be $A[G]$-modules, projective and  of finite type over $A$; denote 
$$V^{D}:= \Hom(V, A(d)).$$
Let $f_{i}\in \pi_{i}:=\Hom( M_{i}, V_{i})$ be  $A[G]$-maps,  suppose we have fixed an identification $V_{2}^{D}\cong V_{1}$;  let $\la\ , \ \ra$ be a perfect pairing $M_{1}\times M_{2}  \to A(d)$, inducing $u_{\la ,\ra} \colon M_{2}^{D}\to M_{1}$. Let $f_{2}^{D}\colon V_{2}^{D}\cong V_{1}\to M_{2}^{D}$ be the dual map, then 
 we define a pairing on $\pi_{1}\times \pi_{2}$ by
\beq\label{pairinggg}
\source{universal-p-adic-gross-zagier-formula}
(f_{1}, f_{2})_{\la, \ra}= f_{1}\circ u_{\la , \ra}(f_{2}^{D}) \in \End_{A[G]}(V_{1}).\eeq
\end{enonce}


\subsubsection{Homological dualities / 1} 
Fix lattices $W^{\circ}$ and $W^{\vee, \circ}$ on any  right algebraic representation of $\G_{*}$ over $L$, and denoted by $\la\ , \ \ra^{W}\colon W\ot W^{\vee}\to L$ the natural invariant pairing. This may not preserve the lattices but it does so up to a bounded denominator which we denote by $p^{-|W|}$.\footnote{With  respect to the model in \eqref{inv pair W}, we have $|W|={\rm ord}_{p}\left({k-2\choose (k-2+l)/2}\right)$ for the representation \eqref{Wwkl} of $\G$.}









We may then consider the  Poincar\'e duality pairings
\beq\label{poincare} \la , \ra_{U_{*},W}\colon H_{d}(\baar{X}_{*, U_{*}}, \cW)\times H_{d}(\baar{X}_{*, U_{*}}, \cW^{\vee})\to H_{0}(\baar{X}_{*, U_{*}}, \cW\otimes \cW^{\vee})\to L(d),
\source{universal-p-adic-gross-zagier-formula}
\eeq
where the second map is induced by $\la, \ra^{W}$ and summation over the connected components of $\baar{X}_{*, U_{*}}$.
  These pairings are integral up to  a bounded denominator $p^{-|W|}$ and  satisfy 
$$\la xT, y\ra_{U_{*}, W}= \la x, yT^{\iota}\ra_{U_{*}, W} $$
for any $T$ in $\cH_{\G_{*}, U_{*}}^{p}$, as well as the projection formula
\beq\label{proj formula}
\source{universal-p-adic-gross-zagier-formula}
\la {\rm p}_{U'_{*}/U_{*}, *}(x), y\ra_{W, U_{*}} = \la x,  {\rm p}_{U'_{*}/U_{*}}^{*}(y)\ra_{W, U_{*}'} 
\eeq
for all pairs of levels $U_{*}'\subset U_{*}$; here ${\rm p}_{U'_{*}/U_{*}}\colon X_{U'_{*}}\to X_{U_{*}}$ is the projection. 


\subsubsection{Homological dualities / 2} We start to promote and modify the Poincar\'e duality pairings. The following lemma is clear.


\begin{lemm} \label{et trace}
\source{universal-p-adic-gross-zagier-formula}
 Let $R$ be a ring, $S$  a finite $R$-algebra,   $M$ a finite $S$-module. 
\begin{enumerate}
\item Suppose that $S$ is \'etale over $R$. Then there is a natural isomorphism
 $$\alpha\colon \Hom_{R}(M, R)\to \Hom_{S}(M, \Hom_{R}(S, R)) \to \Hom_{S}(M, S)$$
 where the first map is $\lambda\mapsto (m\mapsto (s\mapsto \lambda(s m)).$
and the  second one comes from the isomorphism $S\cong \Hom_{R}(S, R)$ induced by the relative trace map.
 \item Suppose that $S=R[T]$ for a finite abelian group $T$, then there is an isomorphism
$\beta\colon \Hom_{R}(M, R)\to \Hom_{R[T]}(M, R[T])$
given by $\lambda\mapsto (m\mapsto \sum_{t}\lambda(tm) [t^{-1}])$.
\end{enumerate}
If $S=R[T]$ is \'etale over $R$ then we have $\alpha(\lambda)=|T|^{-1}\beta(\lambda)$.
\end{lemm}

If $S=R[T]$ for a finite abelian group $T$, one verifies that the   isomorphism of the lemma is given by
$$\la\ ,  \  \ra\mapsto \la\la\ ,\ \ra \ra, \quad \la\la x, y\ra\ra:= \sum_{t\in T} \la x, ty\ra [t^{-1}].$$


We may apply case 2 of the  lemma to  
$$M=M_{\H,V^{p}, r,W}\otimes M_{\H, V^{p},r,W^{\vee}}, %= H_{0}(\baar{Y}_{V^{p},r}, \cW)$, $M_{2}=M_{\H, V^{p},r,W^{\vee}}$, 
\quad R=L,\quad  S= \Lambda_{\H, V^{p},r}:=  \Lambda_{\H, V^{p},r}^{\circ}\ot_{\OO_{L}}L \cong L[{\baar{T}_{\H, 0}/ \baar{T}_{\H, r}}]$$ (with the isomorphism of \eqref{lambda r}). 
We obtain, from the pairings $\la\ , \ra_{V_{r}, W}$, pairings
$$\la \la \ , \ \ra\ra_{V^{p},W, r}\colon M_{\H,V^{p},  W, r}\otimes_{\Lambda_{\H, W, r}} M_{\H, V^{p},  W^{\vee}, r}^{\iota} \to \Lambda_{\H, V^{p}, W, r}\otimes L,$$
and thanks to an easily verified compatibility, a well-defined pairing 
\beq\label{pair EH}\la \la \ , \ \ra\ra_{V^{p},W} \colon M_{\H,V^{p},W}\otimes_{\Lambda_{H}} M_{\H,V^{p}, W^{\vee}}^{\iota} &\to \Lambda_{\H}\otimes L=\OO_{\cE_{\H}}\otimes L\\
\source{universal-p-adic-gross-zagier-formula}
x\otimes y&\mapsto \lim_{r }  \la \la x_{r}, y_{r} \ra\ra_{V^{p},W, r}.
   \eeq


\subsubsection{Automorphic inner products}


Let 
$${\rm v}(U_{*}):= \vol (X_{*, U_{*}}(\bC))$$ where `$\vol$' denotes the volume with respect to the metric deduced from  the hyperbolic metric $dxdy/2\pi y^{2}$ (using the complex uniformisation \eqref{z quot}),
when $\G_{*}=\G$, and  the counting metric, when $\G_{*}=\H$. By \cite[Lemma 3.1]{yzz}, ${\rm v}(U_{*})\in \Q^{\times}$ and, when $d=\dim X_{*}=1$, it equals the degree of the Hodge bundle $L_{U_{*}}$ defined  as in \emph{loc.cit}. We have
\beq \label{deg vol}
\source{universal-p-adic-gross-zagier-formula}
\deg {\rm p}_{U_{*}' ,U_{*} } ={\rm v}(U'_{*})/{\rm v}(U_{*}) = {\rm Z}_{\G_{*}}(\Q)\cap U_{*}\bks U_{*}/U_{*}',
\eeq
where the last equality can be easily seen e.g. from the complex uniformisation \eqref{z quot}.
We set for any $r\geq 1$
\beq\lb{vUpp}{\rm v}(U_{*}^{p}):={{\rm v}(U_{*}^{p}U_{*,0}(p^{r})_{p})\over p^{dr[F:\Q]}}\eeq
 where $U_{*,0}(p^{r})_{p}\subset \G_{*}(\Q_{p})$ is a maximal compact subgroup if $\G_{*}=\H, \H'$, it is the group of those matrices that are upper triangular modulo $p^{r}$ if $\G_{*}=\G$, and it is deduced from those by product and quotient if $\G_{*}=\G\ts\H, (\G\ts\H)'$.  The right hand side of \eqref{vUpp} is independent of $r\geq 1$.


Let $\pi$ be an automorphic representation of $\G_{*}(\A^{\infty})$ of weight $W^{*}$ over $L$, $V_{\pi}$ the corresponding $G_{E_{*}}$-representation. Then we have an isomorphism $V_{\pi^{\vee}}\cong V_{\pi}^{*}(1)$, hence we may use   Construction \ref{constr pairing}
 with $A=L$, $G=G_{E_{*}}$,  $M_{1}=  \H_{d}(\baar{X}_{*,U_{*}}, \cW)$, $M_{2}=  \H_{d}(\baar{X}_{*,U_{*}}, \cW^{\vee})$, $V_{1}=V_{\pi}$, $V_{2}=V_{\pi^{\vee}}$ and the pairings \eqref{poincare}.
  Using \eqref{carayol iso 2}, we obtain 
 $$( \ , \  )_{\pi, U_{*}}:=(\ ,  \ )_{\la, \ra_{U_{*},W}}\colon \pi^{U_{*}}\times \pi^{\vee,U_{*}}\to L.$$ 
One verifies thanks to \eqref{proj formula} and  \eqref{deg vol} that  the pairing 
\beq\label{pair pi}
\source{universal-p-adic-gross-zagier-formula}
(  \ , \ )_{\pi}:= \lim_{U_{*} }\ (\dim W^{} \cdot {\rm v}(U_{*}))^{-1}\cdot  ( \ , \  )_{\pi,{U_{*}}}\colon \pi\times \pi^{\vee}\to L.
\eeq
is well-defined.


When  $\G_{*}=\H$, denoting $\pi=\chi_{\H}$, we may alternatively apply Construction \ref{constr pairing} to $A, M_{1}, M_{2}, V_{1}, V_{2}$ as above and the image of the pairings $\la \la \ ,  \ra\ra_{V^{p},r W} $ under the map $\Lambda_{\H, r, W}\to L$ given by $[t]\mapsto \chi_{\H}(t)$, and denote the resulting pairings on $\chi_{\H}\times \chi_{\H}^{-1}$ by $(\ , \ )_{\la\la, \ra\ra_{\chi_{\H},V^{p},r}}$. As $|\baar{T}_{\H, 0}/\baar{T}_{\H, r}|\cdot {\rm v}(V^{p})={\rm v}(V^{p}V_{p,r})$ by \eqref{deg vol}, we have 
$$( \ ,  \ )_{\chi_{\H}}  ={\rm v}(V^{p})^{-1} (\ , \ )_{\la\la, \ra\ra_{\chi_{\H},V^{p},r}},$$
and in particular the right-hand side is independent of $V^{p}$.







Assume for the rest of this subsection that $\G_{*}=\G,\G\times \H, (\G\times\H)'$. Then we need  a twist  in order to isolate the toric action and to obtain the $\iota$-equivariance  of the  pairings under the action of the $\Up_{p\infty}$-operators. 



Let $\pi=\pi^{\infty}\ot W$ be an ordinary representation of $\G_{*}(\A)$. Using the transformation
$w_{\rm a}^{\ord}$
defined in Proposition \ref{w ord},  we define a pairing
\beq\label{()'} (f_{1} , f_{2} )_{\pi}^{\ord}:= \dim W^{}  \cdot (w_{\rm a}^{\ord}f_{1}, f_{2})_{\pi}
\source{universal-p-adic-gross-zagier-formula}
 \colon \pi^{\ord}\times \pi^{\vee, \ord} \to L.
\eeq
See Lemma  \ref{cor w ord} for its nondegeneracy.

\subsubsection{Homological dualities / 3} Analogously to the previous paragraph, we define a twisted Poincar\'e pairing
\beq\label{lara'}
\source{universal-p-adic-gross-zagier-formula}
 \H_{1}(\baar{X}_{*,U_{*}^{p},r}, \cW)^{\ord}\otimes  \H_{1}(\baar{X}_{*,U_{*}^{p}, r}, \cW^{\vee})^{\ord} &\to L(1)\\
\la x, y\ra_{U_{*}^{p},W, r}^{\ord} &:= \la x, y  w_{\rm a}^{\ord}\ra_{U_{*}^{p}U_{*,p,r},W},
\eeq
of which we will especially consider the restriction to the ordinary parts of homology.



\begin{lemm}
Let $\pi$ be an ordinary representation of $\G_{*}(\A)$, and identify $$\pi^{\ord} = \Hom_{L[G_{E_{*}}]}(\H_{1}(\baar{X}_{*,U_{*}^{p},r},\cW)^{\ord}, V_{\pi})$$ for sufficiently large $r$ similarly to   Proposition \ref{fibreM2}. Then Construction \ref{constr pairing} provides a pairing $(,)_{\la, \ra_{W, r}^{\ord}}$ on $\pi^{\ord}\times \pi^{\vee,\ord}$; it is related to 
 \eqref{()'} by
\beq\label{compare'}
\source{universal-p-adic-gross-zagier-formula}
(,)^{\ord}_{\pi} =   
{ {\rm v}(U_{*})}^{-1}\cdot
  (, )_{\la, \ra^{\ord}_{U*, W, r}}.\eeq%(,)^{\ord}_{\pi} =   
\end{lemm}
\begin{proof} This follows by chasing the definitions.\end{proof}






By applying case 2 of Lemma \ref{et trace} as in \eqref{pair EH}, corrected by a factor $p^{r[F:\Q]}$,\footnote{This factor accounts for the `$K_{0}(p^{r})$'-part of the level.}
 we obtain from \eqref{lara'}
 pairings
\beq\label{larara}
\source{universal-p-adic-gross-zagier-formula}
 \la \la , \ra\ra_{U^{p}_{*}, W, r}\colon \H_{1}(\baar{X}_{*,U_{*}^{p},r}, \cW)\otimes_{\Z_{p}}  \H_{1}(\baar{X}_{*,U_{*}^{p}, r}, \cW^{\vee}) &\to \Lambda_{\G_{*},U_{*}^{p}, r}(1) \\
x\ot y& \mapsto p^{r[F:\Q]} \sum_{t\in \baar{T}_{\G_{*,0}}/\baar{T}_{\G_{*, r}}} \la x, y\ra_{U^{p}_{*}, W, r}^{\ord}
\eeq

\begin{lemm} \lb{Tequiv}
The parings \eqref{larara} satisfy $\la \la x_{r}T, y_{r}\ra\ra_{W, r}=\la \la x_{r}, y_{r}T^{\iota}\ra\ra_{U^{p}_{*},W, r}$ for all $T\in \cH_{\G, r}^{\ord}$ and all $x_{r}\in  H_{d}(\baar{X}_{*,U_{*}^{p},r}, \cW)$, $ y_{r}\in H_{d}(\baar{X}_{*,U_{*}^{p},r}, \cW^{\vee})$. 

For $z\in M_{\G_{*}, \cW}$, denote by   $z_{r}$ its image in $M_{\G_{*}, \cW, r}:=H_{d}(\baar{X}_{*,U_{*}^{p},r}, \cW)^{\ord}$. The pairing
\beq\label{pair lam}
\source{universal-p-adic-gross-zagier-formula}
\la \la \ , \ \ra\ra_{\Lambda, U^{p}_{*},W} \colon  M_{\G_{*},U_{*}^{p},\cW}\otimes_{\cH_{\G}^{\ord}}  M_{\G_{*},U_{*}^{p},\cW^{\vee}}^{\iota} & \to \Lambda_{\G_{*}, U_{*}^{p}}(d)\otimes L\\
\la\la x , y^{\iota} \ra\ra_{\Lambda,U^{p}_{*}, W}&:=\lim_{r}\la\la x_{r}, y^{\iota}_{r}\ra\ra_{U^{p}_{*},W, r}
\eeq
is well-defined.
\end{lemm}
The above construction is a  minor variation on the one of  \cite[\S~ 2.2.4]{fouquet}, to which we refer for the proof of the lemma.
As usual, when $W=\Q_{p}$ we shall omit it from the notation.


\begin{lemm}\label{pair indep} The diagram
\source{universal-p-adic-gross-zagier-formula}
\begin{equation*}
\xymatrix{
M_{\G_{*},U_{*}^{p}, \cW}\otimes_{\cH_{\G_{*}}^{\circ}} M_{\G_{*},U_{*}^{p}, \cW^{\vee}}^{\iota}\ar[r]^{{\la\la\, ,\,  \ra\ra}_{\Lambda,U^{p}_{*}, W}} &\Lambda_{\G_{*}, U^{p}_{*}}(d)\otimes L\\
M_{\G_{*}, U_{*}^{p}}\otimes_{\cH_{\G_{*}}^{\circ}} M_{\G_{*}, U_{*}^{p}}^{\iota}\ar[r]^{\la\la\, , \, \ra\ra_{\Lambda, U^{p}_{*}}} \ar[u]^{j_{W}\otimes j_{W^{\vee}}} &\Lambda_{\G_{*}, U^{p}_{*}}(d)\otimes L\ar[u]^{\cong},
}
\end{equation*}
where the left vertical map comes from Proposition \ref{free indep}.2 and the right vertical map is  $[t]\mapsto \sigma_{W}^{-1}(t)[t]$, is commutative.
\end{lemm}
\begin{proof} For simplicity we write down the proof for the group $\G_{*}=\G$ and we drop the subscripts $U^{p}$.
Poincar\'e duality and  the  pairings $ \la\ , \ \ra^{W} $  preserve integral structures up to $p^{-|W|}$.  Then  by construction it suffices to show the identity
$$\la j_{W, r}(x), j_{W^{\vee},r}(y)\ra_{W}^{\ord} \equiv  \la x, y \ra^{\ord} \pmod{p^{r-|W|}\OO_{L}}$$
for all $r\geq 1$ and  $x$, $y$ in $ H_{d}(\baar{X}_{ r},\Z_{p})$. 

By definition in \eqref{jWr}, we have
$$j_{W, r}(x)=x\ot \zeta\ot \zeta^{*}$$ where $\zeta_{r}\in W^{\circ, N_{0}}/p^{r}$ and $\zeta_{r}^{*}\in W^{\circ}_{N_{0}}/p^{r}$ are elements pairing to $1$; we denote by $\zeta_{r}^{\vee}$, $\zeta_{r}^{\vee, *}$ the analogous elements for $j_{W^{\vee}, r}$. 
 Then we need to show that 
$$\la( x\ot \zeta_{r}\ot \zeta_{r}^{*}) w_{\rm a}^{\ord}, y\ot \zeta_{r}^{\vee}\ot \zeta_{r}^{\vee*} \ra= \la x, y\ra.$$
By the definition of $w_{\rm a}^{\ord}$ in Proposition \ref{w ord}, this reduces to the identity
$$\la \zeta_{r} w_{0}, \zeta_{r}^{\vee}\ra \cdot \la \zeta_{r}^{*} w_{0}, \zeta_{r}^{\vee, *}\ra = \la \zeta_{r} , \zeta_{r}^{*}\ra\cdot \la \zeta_{r}^{\vee} , \zeta_{r}^{\vee,*}\ra=1, $$
which can be  immediately  verified using an explicit model for the pairing such as given in \eqref{inv pair W}. 
\end{proof}


\subsubsection{Dualities over Hida families} 

Let $\X$ be an irreducible component of $\cE_{K^{p}}^{\ord}$. By Proposition \ref{etale}, the map $\cE^{\ord}_{K^{p}}\to \Spec \Lambda_{\Q_{p}}$ is \'etale  in a neighbourhood $\X'$ of $\X^{\cl}$, hence we may  apply case 1 of  Lemma \ref{et trace} to deduce  from \eqref{pair lam} a pairing 
\beq
\label{pair E}
\source{universal-p-adic-gross-zagier-formula}
\la \la \ , \ \ra\ra_{K^{p}{}^{\ord}} \colon \cM_{K^{p}{}^{\ord}}\otimes_{\OO_{\X'}} \cM_{K^{p}{}'}^{\iota} \to \OO_{\X'}(1).\eeq





We summarise the situation.

\begin{prop}[Duality]
\source{universal-p-adic-gross-zagier-formula}
\notready\label{duality}
\source{universal-p-adic-gross-zagier-formula}
Let $\X^{(5)}\supset \X^{\cl}$ be the intersection of the subset $\X^{(4)} $ of Theorem \ref{LGC} with the locus where the map $\X\to \Spec\Lambda_{\Q_{p}} $ is \'etale. There exist
\begin{itemize}
\item  a perfect, $G_{E}$-equivariant, skew-hermtian pairing 
\begin{align}
\cM_{K^{p}{}'}^{{H_{\Sg}'}}\otimes_{\OO_{\X^{(5)}}} \cM_{K^{p}{}'}^{{H_{\Sg}'},\iota}\to \OO_{\X^{(5)}}(1).\label{iso 1}
\source{universal-p-adic-gross-zagier-formula}
\end{align}
induced from \eqref{pair E};
\item  a perfect, $G_{E}$-equivariant, skew-hermtian pairing 
\beq 
\cV \ot_{\X^{(5)}} \cV^{\iota} \to \OO_{\X^{(5)}}(1).\label{iso 2}.
\source{universal-p-adic-gross-zagier-formula}
\eeq
\item a perfect pairing 
\begin{align}\label{iso5}
\source{universal-p-adic-gross-zagier-formula}
((\ , \  )):= {\rm v}(K^{p}{}')^{-1}\cdot ( \ , \  )_{\la \la, \ra\ra_{K^{p}{}'}} \colon \Pi_{{H_{\Sg}'}}^{K^{p}{}'}\otimes_{\OO_{\X^{(5)}}} (\Pi_{{H_{\Sg}'}}^{K^{p}{}'})^{\iota}  \to \OO_{\X^{(5)}},\end{align}
where $(\ , )_{\la \la, \ra\ra_{K^{p}{}'}}$ is deduced from \eqref{iso 1}, \eqref{iso 2} and the isomorphism of  Proposition \ref{cofinal}.2 via Construction \ref{constr pairing}. 
\end{itemize}
\end{prop}
\begin{proof} 
Observe that the natural map
$$ (\cM_{K^{p}{}'}^{{H_{\Sg}'}})^{*}\to ((\cM_{K^{p}{}'})_{{H_{\Sg}'}})^{*}  =(\cM_{K^{p}{}'}^{*})^{{H_{\Sg}'}}$$
(where ${}^{*}$ denotes $\OO_{\X^{(5)}}$-dual) is an isomorphism; as \eqref{pair E} is equivariant for the action of the full Hecke algebra, this implies that its restriction \eqref{iso 1} is perfect.  It is skew-hermitian by Lemma \ref{Tequiv} and the fact that the Poincar\'e pairing  \eqref{poincare}, when $W$ is trivial, is skew-symmetric. 


 We find the pairing \eqref{iso 2} by specialising \eqref{iso 1} to   $K^{p}{}'=K^{p}$, and the pairing \eqref{iso5} as described.
\end{proof}




\subsubsection{Specialisations} We describe the specialisation of the pairing $((, ))$ just constructed. 

For each algebraic representation $W $ of $(\G\times \H)$, denote by
$$\X_{r}^{\cl,W}:=\X\cap \cE_{K^{p},r}^{\cl , W},$$
the set of classical points of weight $W$ (omitted from the notation when $W=\Q_{p}$) and level $r$. Denote by a subscript `$W,r$' the pullbacks of sheaves or global sections from $\X^{(5)}$ to $\X_{r}^{\cl, W}$ (which is a finite \'etale scheme  over $\Q_{p}$). We let 
\begin{gather*}
V_{W,r}:=\Gamma(\X_{r}^{\cl, W}, \cV_{W, r}), \qquad
\Pi_{{H_{\Sg}'}, W, r}^{K^{p}{}', \ord}:=\Gamma(X_{r}^{\cl, W}, \Pi_{{H_{\Sg}'}}^{K^{p}{}', \ord}), \\
M_{K^{p}{}', W, r}^{{H_{\Sg}'}}:= \Gamma(\X_{r}^{\cl, W},  \cM_{K^{p}{}'}^{{H_{\Sg}'}})=  H_{1}(\baar{Z}_{K^{p}{}',r} , \cW)^{\ord,{H_{\Sg}'}},
\end{gather*}
where the last equality is by Proposition \ref{free indep}.3. 
We  denote
$$((\ , \ ))_{W, r}:={\rm v}(K^{p}{}')^{-1} \cdot (\ , \ )_{\la \la, \ra\ra_{K^{p}{}', W, r}}\colon \Pi_{{H_{\Sg}'}, W, r}^{K^{p}{}', \ord}\times \Pi_{{H_{\Sg}'}, W, r}^{K^{p}{}', \ord, \iota}\to \OO(\X_{r}^{\cl, W}).$$




\begin{prop}
\source{universal-p-adic-gross-zagier-formula}
\notready \label{abovco} % and let $M_{K^{p}{}' , r}$, respectively $V_{r}$, be the spaces of global sections of $\cM_{K^{p}{}', r}$, respectively $\cV_{r}$. 
\source{universal-p-adic-gross-zagier-formula}
Let $f_{1}$, respectively $f_{2}$ be global\footnote{The same statements hold with some extra notational burden if $f_{1}, f_{2}$ are only defined over an open subset of $\X^{(5)}$.} sections of $\Pi^{K^{p}{}', \ord}_{{H_{\Sg}'}}=\underline{\Hom}_{\OO_{\X^{(5)}}[G_{F, E}]}(\cM_{K^{p}{}'}^{{H_{\Sg}'}},\cV) $, respectively $(\Pi^{K^{p}{}', \ord}_{{H_{\Sg}'}})^{\iota}$. 
Let $$f_{1, W, r}\colon  M_{K^{p}{}', W, r}^{{H_{\Sg}'}} 
 \to V_{W,r}, \qquad f_{2, W, r}\colon 
( M_{K^{p}{}', W, r}^{{H_{\Sg}'}})^{\iota}\to V_{W,r}^{\iota}
$$ be $\OO_{\X^{\cl, W}_{r}}[G_{F, E}]$-linear maps. 

Let $x\in \X^{\cl, W}_{r}$, $\pi:= \pi(x) $ and let 
$$ f_{1, x}\colon  \H_{1}(\baar{Z}_{K^{p}{}',r} , \cW)\to V_{\pi},\qquad f_{2, x}\colon \H_{1}(\baar{Z}_{K^{p}{}',r} , \cW^{\vee})\to V_{\pi^{\vee}}$$  
be $\Q_{p}(x)[G_{F, E}]$-linear maps.

The following hold.
\begin{enumerate}
\item Suppose that for $i=1,2$,   the map $f_{i, W, r}$ arises as the specialisation of $f_{i}$. Then 
 $${((f_{1}, f_{2}))}|_{\X_{r}^{\cl, W}}= ((f_{1, W, r}, f_{2, W, r}))_{{W,r}}\qquad \text{in } \OO(\X_{r}^{\cl, W}).$$
 \item Suppose that  for $i=1, 2$, the map $f_{i, x}$ factors through the projection
 $${\rm p}^{?}\colon  \H_{1}(\baar{Z}_{K^{p}{}',r} , \cW^{?})\to  \H_{1}(\baar{Z}_{K^{p}{}',r} , \cW^{?})^{\ord}_{{H_{\Sg}'}} \cong H_{1}(\baar{Z}_{K^{p}{}',r} , \cW)^{\ord,{H_{\Sg}'}} = M_{K^{p}{}', W, r}^{{H_{\Sg}'}},$$
 where $?=\emptyset$ if $i=1$, $?=\vee$ if $i=2$; 
 and that $f_{i,x}$ coincides with the specialisation of $f_{i, W, r}$ at $x$. Then 
 \beq\label{eq specialise}
\source{universal-p-adic-gross-zagier-formula}
((f_{1}, f_{2}))_{W, r }(x)
= (f_{1,x}, f_{2,x})^{\ord}_{\pi} =\dim W\cdot (w_{\rm a}^{\ord} f_{1}, f_{2})_{\pi}
 \qquad \text{in } \Q_{p}(x).
\eeq
\end{enumerate}
\end{prop}

\begin{proof}
We simplify the notation by omitting the superscripts $ {{H_{\Sg}'}}$ and subscripts $K^{p}{}'$; moreover we ignore the normalisations ${\rm v}(K^{p}{}')^{-1}$ that are present in all of the pairings to be compared.

 
Part 1 follows from the definition \eqref{pair lam} if $W=\Q_{p}$, and similarly we can also identify $((\ ,  \ ))_{W,r}$ with  the restriction to $\X_{r}^{\cl, W}$ of the pairing on functions on $M_{W}$ deduced from $\la\la\ , \ \ra\ra_{W}$ via 
Construction \ref{constr pairing}. By Lemma \ref{pair indep} this implies that the desired statement holds for all $W$.


For Part 2, let $H_{W,r}:=\OO(\X_{r}^{\cl , W})$.  
First  notice that, by the construction of case 1 of Lemma \ref{et trace}, the diagram  of $H_{W, r}$-modules
\beqq
\xymatrix{
 \Hom_{H_{W,r}}(M_{W,r}^{\iota}, H_{W,r}(1))_{\baar{\Q}_{p}}\ar[d]\ar[r]^{\ \ \ \  \ \  \ u_{\la\la, \ra\ra_{W, r}}}  & M_{W,r}\ar[d]\\
  \Hom(M_{W, r, \Lambda_{W, r}}^{\iota}, \Lambda_{W, r}(1))\ar[r]^{\ \   \ \  \ \ \ u_{\la \la, \ra\ra_{\Lambda, W, r}}}  & M_{W,r}  
} 
\eeqq
is commutative. 

On the other hand, let $x\in \X_{\cl}^{W, r}$ and let $\alpha_{x}$ be the associated character of $T^{+}$. By  definition in \eqref{larara},
the pairing $\la  \la, \ra\ra_{\Lambda, W}$ specialises, on $M_{W,r|x}\otimes M_{W, r|x}^{\iota}$, to 
$$(x, y)\mapsto \sum_{t\in \baar{T}_{0}/\baar{T}_{r}} \la x, ty\ra_{U_{*}^{p},W, r}^{\ord} [t^{-1}](x)= \sum_{t\in \baar{T}_{0}/\baar{T}_{r}} \alpha_{x}(t) \la x, y\ra_{U_{*}^{p},W, r}^{\ord} \alpha^{-1}_{x}(t) =p^{r[F:\Q]} |\baar{T}_{0}/\baar{T}_{r}|\cdot \la x, y\ra_{U_{*}^{p},W, r}^{\ord}.$$
It follows that $u_{\la\la, \ra\ra_{W, r}}$ specialises at $x$ to
 $p^{-r[F:\Q]}|\baar{T}_{0}/\baar{T}_{r}|^{-1}u_{\la, \ra'_{W,r}}$, hence that the specialisation of $((, ))(x)={\rm v}({K^{p}{}'})^{-1}(,)_{\la\la,\ra\ra_{W,r}}(x)$ is
$${  (,)_{\la,\ra_{W,r}^{\ord}} 
\over p^{r[F:\Q]} |\baar{T}_{0}/\baar{T}_{r}|  {\rm v}({K^{p}{}'})}
 =
{p^{r[F:\Q]}
\cdot
{ {\rm v}(K^{p}{}' K_{1}^{1}(p^{r})_{p})}  (, )_{\pi}^{\ord}
\over {p^{r[F:\Q]} [{\rm v}(K^{p}{}' K_{1}^{1}(p^{r})_{p })/ {\rm v}(K^{p}{}' K_{0}(p)^{r}_{p }) ]   {\rm v}(K^{p}{}' K_{0}(p^{r})_{p})}}
= (, )_{\pi}^{\ord},$$
where we have used $|\baar{T}_{0}/\baar{T}_{r}|= {\rm v}(K^{p}{}' K_{0}(p^{r})_{p }) / {\rm v}(K^{p}{}' K_{1}^{1}(p)^{r}_{p }  )$ (by \eqref{deg vol}) and \eqref{compare'}. 

This establishes the first equality of \eqref{eq specialise}; the second one is just a reminder of \eqref{()'}.
\end{proof}



\subsection{Local toric pairings} \lb{sec 42}
Let $F$ be a non-archimedean local field, $E$ a quadratic \'etale algebra over $F$ with associated character $\eta\colon F^{\times}\to\{\pm 1\}$, $B$ a quaternion algebra over $F$,  $G=B^{\times}$, $H=E^{\times}$, $H'=H/F^{\times}$, and suppose given an embedding $ H\into G$ 

\subsubsection{Definition of the pairing}
Let $\pi$ be a smooth irreducible representation of $G$ over a finite extension  $ L$  of $\Q_{p}$, with a central character $\omega\colon F^{\times}\to L^{\times}$. Let $\chi\colon E^{\times}\to L^{\times}$ a character such that $\chi|_{F^{\times}}\cdot \omega=\one$. We identify $\chi$ with a representation  $L\chi$ of $E^{\times} $ on $L$, and when more precision is needed we denote by $e_{\chi}$ the basis element corresponding to the character $\chi$ in $L\chi$.  Let  $\Pi:=\pi\otimes\chi$, a representation of $(G\ts H)'=(G\times H)/F^{\times}$ over $L$. We assume that $\pi$ is essentially unitarisable, that is  that for any embedding $\iota\colon L\into \bC$,  a twist of $\iota\pi$ is isomorphic to the space of smooth vectors of a unitary representations. (This holds automatically if $\pi $ arises as the local component of a cuspidal automorphic representation over $L$.)
Let $\pi^{\vee}$ be the smooth dual, $\Pi^{\vee}:=\pi^{\vee}\otimes \chi^{-1}$

Assume from now on  that the modified local sign  $\eps(\Pi)=\eqref{eps v}$ equals $+1$.
 Then, by the result of Tunnell and Saito mentioned in the introduction, 
the space
$$\Pi^{*,H'}:={\Hom}_{H'} (\Pi, L).$$
has dimension   $1$ over $L$. Moreover the choices of an invariant pairing $(\ , \ )$ on $\Pi \otimes \Pi^{\vee}$  and a Haar measure $dt$ on $\H'$ 
give a   generator  
$$Q=Q_{(\ , \ ), dt}\in \Pi^{*,H'}\otimes_{L} (\Pi^{\vee})^{*,H'}$$
defined by the absolutely convergent integral
\begin{gather}\label{Qpair}
\source{universal-p-adic-gross-zagier-formula}
 Q_{(, )}(f_{1}, f_{2}):=
 \calL(V_{v}, 0)^{-1}
 \cdot
\iota^{-1} \int_{E^{\times}/F^{\times}} (\iota\Pi(t) f_{1}, \iota f_{2})\, dt ;
 \end{gather}
 for any  $\iota\colon L\into \bC$; here $\calL(V_{v}, 0)=\eqref{calLv}$. 
 
 Recall also from the introduction  \eqref{Qv intro} that 
\beq\label{Qv sec4}
\source{universal-p-adic-gross-zagier-formula}
 Q_{dt}\left({f_{1}\ot f_{2}\over f_{3}\ot f_{4} }\right) := {Q_{( , ), dt}(f_{1}, f_{2})\over ( f_{3}, f_{4})},
\eeq
is independent of $(\ ,\ )$ whenever it is defined.

 
 
 
We study the pairing, or some of its variations, in a few different contexts.

 

 
\subsubsection{Interpretation in the case $E=F \oplus F$} In this case   $G=\GL_{2}(F)$, and  the  integral \eqref{Qpair} has an interpretation as product of zeta integrals.  Let $\cK(\pi)$ and $\cK(\pi^{\vee})$ be Kirillov models over $L$ as in \cite[\S~ 2.3]{dd-pyzz}.
By \cite[Lemma 2.3.2]{dd-pyzz}, the $L$-line of invariant pairings on $\cK(\pi)\times\cK(\pi^{\vee})$  is generated 
by an element  $(\ , \ )$ such that, for each $\iota\colon L\into \bC$, 
we have 
\beq \label{kir pair}
\source{universal-p-adic-gross-zagier-formula}
\iota ( f, f^{\vee})={\zeta(2)\over L(1, \pi\times \pi^{\vee})}\cdot \int_{F^{\times}} \iota f(y) \iota f^{\vee}(y) d^{\times }y,\eeq
where
 the integral is absolutely convergent (as $\iota\pi$ is essentially unitarisable) and $d^{\ts}y$ is any $L$-valued Haar measure. Identify $E^{\times}$ with the diagonal torus in $\GL_{2}(F)$ and write $\chi=(\chi_{1}, \chi_{2})$ according to the decomposition $E=F\oplus F$; noting that $\chi_{2}=\omega^{-1}\chi_{1}$ and $\pi=\pi^{\vee}\otimes \omega^{-1}$, we indentify   $Q_{(, ), dt}$ with
\begin{gather}\label{Qv reg}
\source{universal-p-adic-gross-zagier-formula}
\iota Q_{(, ), dt}(f\otimes e_{\chi}, f^{\vee}\otimes e_{\chi^{-1}})\doteq L(1/2, \iota \pi_{E}\otimes \iota\chi)^{-1} \int_{E^{\times}/F^{\times}} \iota \chi(t) (\iota\pi(t) f,  \iota f^{\vee} )|t|^{s}\, dt|_{s=0}\\
\nonumber
=
 L(1/2, \iota \pi\otimes \iota\chi_{1})^{-1} \int_{F^{\times}} \iota f^{\vee}(t)\iota\chi_{1}(t) |t|^{s}d^{\times }t|_{s=0}
\cdot
L(1/2, \iota \pi\otimes \iota\chi_{1})^{-1} \int_{F^{\times}} \iota f^{\vee}{}(y)\iota\chi_{1}(y) |y|^{s}d^{\times }y|_{s=0}\\
\nonumber
= (L(1/2,\iota \pi\otimes \iota \chi_{1})^{-1} \cdot I(\iota f, \iota\chi_{1}, 1/2))\cdot  (L(1/2, \pi\otimes \iota\chi_{1})^{-1} \cdot I(\iota f^{\vee} , \iota\chi_{1} , 1/2)),
\end{gather}
where $I(\cdot , \cdot , 1/2)$ is the zeta integral of \cite[\S~5.2]{LLC} 
 for $\GL_{2}(F)\times \GL_{1}(F)$, and $\doteq$ denotes an equality up to constants in $L^{\ts}$ depending on the choices of measures.

\subsubsection{Special line in the unramified case} We study the first one in a short list of special cases in which there are `canonical' lines in $\Pi$, $\Pi^{\vee}$, on which  the value of the pairings $Q $ can be explicitly computed. 
\begin{lemm}[{\cite[Lemme 14]{wald}}]\label{unrQ} Suppose that $B$ is split, $E/F$ is unramified, and both $\pi$ and $\chi$ are unramified. Let $K\subset (G\times E^{\times})/F^{\times}$ be   a maximal compact subgroup. Then 
\source{universal-p-adic-gross-zagier-formula}
$$Q_{( \ ,\  ), dt}(v, w)=\vol(\OO_{E}^{\ts}/\OO_{F}^{\ts}, dt)\cdot  (v, w)$$
 for all $v$, respectively $w$, in the lines $\Pi^{K}$, respectively $(\Pi^{\vee})^{K}$.
\end{lemm}

\subsubsection{Special line when   $B$ is nonsplit} Suppose now that $B$ is nonsplit and that $\Pi$ is an irreducible representation of $(G\times E^{\times})/F^{\times}$ as above. Note that $\Pi$ is finite-dimensional and $H'$ is compact, so that $\Pi^{\vee}=\Pi^{*}$ and  the natural maps $\Pi^{H'}\to \Pi_{H'}$ ($=H'$-coinvariants) and $\Pi^{*,H'} \to (\Pi^{H'})^{*}$ are isomorphisms. Moreover the non-degenerate pairing $(\ , \ )$ restricts to a non-degenerate pairing on $\Pi_{H'}\otimes \Pi^{\vee}_{H'}$.
Then we may compare the restrictions of the pairings $Q_{(\ ,  \ )}$   of $(\ , \ )$  to the line $\Pi^{H'}\otimes \Pi^{\vee,H'}$.
\begin{lemm}\label{ramifQ} In the situation of the previous paragraph, we have 
\source{universal-p-adic-gross-zagier-formula}
\beqq
Q_{(\ , \ ), dt}= \calL(V_{(\pi , \chi), v}, 0)^{-1}
 \cdot \vol(E^{\times}/F^{\times}, dt)\cdot (\ , \ )
\eeqq
as elements  of $(\Pi^{H'})^{*}\otimes (\Pi^{\vee,H'})^{*}$. 
\end{lemm}
\begin{proof} 
  This follows from the definition in \eqref{Qpair}, since in this case the integration over the compact set $E^{\times}/F^{\times}$ converges.
\end{proof}

\subsection{Ordinary toric pairings} \lb{sec 43}
We define a variant  for ordinary forms of the pairing $Q$.



\subsubsection{Definition of the ordinary paring} Let $\Pi=\pi\otimes \chi$ be an ordinary automorphic representation of $(\G\times \H)'(\A^{})$ over $L$.
When referring to  local objects considered in the previous paragraphs or products thereof, we append subscripts as appropriate. 


For each $v\vert p$, let $$\mu_{v}^{+}\colon E_{v}^{\ts}\to L^{\ts}$$ be the character by which $E_{v}^{\ts}$ (or equivalently $\prod_{w\vert v} G_{E_{w}}^{\rm ab}$) acts on $V_{\pi,v}^{+} \ot\chi_{v}$, and let ${\rm j}_{v}\in E_{v}$ be the purely imaginary element fixed in \eqref{jv}. Define 
$$\mu^{+}({\rm j}):=\prod_{v\vert p} \mu_{v}^{+}({\rm j}_{v}).$$
For measures $dt_{v}=dt_{v, p}$ on $H'_{v}$, $dt_{v,\infty}$ on $H'_{v, \infty}$ (the latter a merely   formal notion as in the introduction), define
\beq\lb{vol circ} {\vol^{\circ}(H'_{v}, dt_{v})} &:=
 {\vol(\OO_{E, v}^{\ts}/\OO_{F, v}^{\ts}, dt_{v})\over e_{v}
   L(1, \eta_{v})^{-1}},
  \qquad { \vol^{\circ}(H'_{v, \infty}, dt_{v, \infty})} := {\vol (H'_{v, \infty}, dt_{v, \infty}) \over 2^{[F_{v}:\Q_{p}]}}, \\
   {\vol^{\circ}(H'_{p\infty}, dt_{p\infty})}&:= \prod_{v\vert p } {\vol^{\circ}(H'_{v}, dt_{v})} \cdot { \vol^{\circ}(H'_{v, \infty}, dt_{v, \infty})}.
\eeq
The denominators in the right-hand sides are the volumes of $\vol(\OO_{E, v}^{\ts}/\OO_{F, v}^{\ts})$, respectively $\bC^{\ts}/\R^{\ts}$, for the ratio of  (rational normalisations of) selfdual measures, cf. \cite[\S~ 1.6.2]{yzz} and the proof of Proposition \ref{toric period}.


\begin{defi}
\source{universal-p-adic-gross-zagier-formula}
\notready\lb{Q ord def}
Let $dt=dt^{p\infty}dt_{p\infty}$ be a decomposition of the ad\`elic measure $dt$ specified in \eqref{dt norm}. 
Then we define:
\begin{itemize}
\item
 for each $f_{1, p \infty}, f_{3,p \infty}\in \Pi_{p \infty}^{\ord}$, $f_{2, p \infty}, f_{4,p \infty}\in \Pi_{p \infty}^{\vee, \ord}$ with $f_{3,p \infty}\ot f_{4,p \infty}^{\ord}\neq 0$, 
\beq 
\label{Q orddd loc} 
\source{universal-p-adic-gross-zagier-formula}
Q^{\ord}_{dt_{p\infty}}\left( {f_{1,p\infty}\ot f_{2,p\infty}\over f_{3,p\infty}\ot f_{4,p\infty}}\right) :=\mu^{+}({\rm j})^{} 
 {\vol^{\circ}(H'_{p\infty}, dt_{p\infty})} \cdot {f_{1, p\infty}\ot f_{2, p\infty}\over f_{3, p\infty}\ot f_{4, p\infty}}.
\eeq
\item
 for each $f_{1}, f_{3}\in \Pi^{\ord}$, $f_{2}, f_{4}\in \Pi^{\vee, \ord}$ with $(f_{3}, f_{4})^{\ord}\neq 0$, 
\beq 
\lb{Q orddd} 
Q^{\ord} \left( {f_{1}\ot f_{2}\over f_{3}\ot f_{4}}\right)  
:=Q^{p\infty}_{dt^{p\infty}}\left( {f_{1}^{p\infty}\ot f_{2}^{p\infty}\over f_{3}^{p\infty}\ot f_{4}^{p\infty}}\right)
\cdot 
Q^{\ord}_{dt_{p\infty}}\left( {f_{1,p\infty}\ot f_{2,p\infty}\over f_{3,p\infty}\ot f_{4,p\infty}}\right).
\eeq
\end{itemize}
\end{defi}
The normalisation at $p\infty$ is justified by the clean formula of Proposition \ref{compare Qs} below.
\begin{rema}
\source{universal-p-adic-gross-zagier-formula}
\notready \lb{Qord nonzero}Suppose  that $\Pi$ is locally distinguished, so that as explained in the introduction the functional $Q_{dt}$ is nonzero. Then the functional $Q^{\ord}_{dt}$ is also nonzero.
\end{rema}


\subsubsection{Decomposition}
Fix a decomposition $dt= \prod_{v\nmid p \infty } dt_{v} dt_{p\infty}$ such that for all but finitely many $v$, $\vol(\OO_{E,v}^{\ts}/\OO_{F, v}^{\ts})=1$. 
Let $\Sg'$ be a finite set of finite  places of $F$  disjoint from $\Sg $ and $S_{p}$ and containing the other places of ramification  of $\Pi$, and those such that $\vol(\OO_{E,v}^{\ts}/\OO_{F, v}^{\ts})\neq 1$. Let $K^{p}\subset (\G\times \H)'(\A^{p\infty})$ be an open compact subgroup that is maximal away from $S:=\Sg\cup \Sg'$ and such that $\Pi_{v}^{K_{v}}=\Pi_{v}$ for $v\in \Sigma$. 
\begin{lemm} For all $f_{1}, f_{3} \in \Pi^{K^{p}, \ord}_{H_{\Sg}'}$, $f_{2}, f_{4} \in \Pi^{\vee, K^{p}, \ord}_{H_{\Sg}'}$ with $(f_{3}, f_{4})^{\ord}\neq 0$, we have
\beq\label{dec Qord} Q^{\ord}\left(   {f_{1} \ot f_{2}\over  f_{3}\ot f_{4} } \right)
\source{universal-p-adic-gross-zagier-formula}
&=
\prod_{v\in \Sg'}Q_{v, dt_{v}}\left({  f_{1, v} \ot  f_{2, v}\over f_{3,v}\ot f_{4, v}}\right)
 \cdot
 \prod_{v\in \Sg} \vol(E_{v}^{\times}/F_{v}^{\times}, dt)   \calL(V_{(\pi , \chi), v}, 0)^{-1}   {  f_{1, v} \ot  f_{2, v}\over f_{3,v}\ot f_{4, v}}\\
&\cdot 
  {f^{Sp\infty}_{1}\otimes f_{2}^{Sp\infty}\over f^{Sp\infty}_{3}\otimes f^{Sp\infty}_{4}} \cdot Q^{\ord}_{p\infty,dt_{p\infty}}\left( {f_{1,p\infty}\ot f_{2,p\infty}\over f_{3,p\infty}\ot f_{4,p\infty}}\right).
 \eeq
\end{lemm}

\begin{proof} This follows from the definitions and the results of \S~\ref{sec 42}. 
\end{proof}
\subsubsection{Relation to between the toric pairing and its ordinary variant} We  gather the conclusion of the computations from the appendix. 



\begin{prop}
\source{universal-p-adic-gross-zagier-formula}
\notready \lb{compare Qs} Let $\Pi=\pi\ot \chi= \Pi^{\infty} \ot W$ be an ordinary representation of $(\G\ts \H)'(\A)$.  
Let  $w_{\rm a}^{\ord}$ and $\tord$ be the operators defined in Propositions \ref{w ord} and \ref{gHo}. Let $e_{p}(V_{(\pi, \chi)})=  \eqref{epinf L}$ be the interpolation factor of the $p$-adic $L$-function.  For all  $f_{1}$, $f_{3}\in \Pi^{\ord}$, $f_{2}$, $f_{4}\in \Pi^{\ord}$ with $(f_{3}, f_{4})^{\ord}\neq0$, we have
\beqq Q^{}\left(   {\tord(f_{1}) \ot \tord (f_{2})\over  w_{\rm a}^{\ord}(f_{3})\ot f_{4} } \right) =   e_{p}(V_{(\pi, \chi)})\cdot \dim W \cdot  Q^{\ord}\left(   {f_{1} \ot f_{2}\over  f_{3}\ot f_{4} } \right). \eeqq
\end{prop}
\begin{proof}
There is a   decomposition $Q_{p\infty , dt_{p\infty}}^{\ord} = \prod_{v\vert p}Q_{v,dt_{v}}^{\ord}\cdot  \prod_{v\vert p}Q_{v, \infty, dt_{v, \infty}}^{\ord}$, whose terms are defined in \S\S~\ref{A3}-\ref{A4}. The only point worth stressing is  that if $\mu_{v}^{+}$, respectively $\mu_{v, \infty}^{+}$ is the character defined in \S~ \ref{A3mu},\footnote{Note that despite the similar notation, the character $\mu_{v}$ is defined using the Weil--Deligne  representations rather than the continuous Galois representations.} respectively \S~\ref{A4mu}, then the decomposition $\mu^{+}=\mu^{+, {\rm sm}}\mu^{+, {\rm alg}}$  of $\mu^{+}$ into a product of a smooth and an algebraic character is given by 
$\mu^{+, {\rm sm}}=\prod_{v\vert p }\mu_{v}^{+}$, $\mu^{+, {\rm alg}}= \prod_{v\vert p}\mu_{v, \infty}^{+}$. 

Then the result follows from Propositions \ref{toric period} and \ref{compare toric inf}.
\end{proof}



 


\subsection{Interpolation of the  toric pairings}  \lb{sec 44} We interpolate the pairings $Q^{\ord}_{dt}$ along Hida families
\subsubsection{Interpolation of  the local pairings} We use the same notation $F, E$ of \S~\ref{sec 42}.

\begin{lemm}\label{interp Lv} Let $\X$ be a scheme over $\Q$ and let   $r'=(r, N)$ be a Weil--Deligne representation  of $W_{F}$
\source{universal-p-adic-gross-zagier-formula}
on a rank-$2$ locally free sheaf over $\X$.  Suppose that $\X$ contains a  dense subset $\X^{\cl}$ such that $r'_{x}$ is pure for all $x\in \X^{\cl}$. Let $\ad(r')$ be the rank-$3$ adjoint representation. Then there exist an open subset $\X''\subset \X$ containing $\X^{\cl }$ and functions 
$$L(0, r')^{-1}, \qquad L(1, r', \ad)\in \OO(\X'')$$
such that for every $x\in \X'$ we have $L(0, r')^{-1}(x)=L(0, r'_{x})^{-1}$ and $L(1, r', \ad)(x)=L(1, \ad(r'_{x}))$.
\end{lemm}
\begin{proof}  
By  \cite[\S~5.1]{LLC},
 there exist an open set $\X'''\subset \X$ containing $\X^{\cl}$ and functions
 $L(0, r')^{-1}$, respectively $L(1, r', \ad)^{-1}$, in $\OO(\X''')$ interpolating $L(0, r'_{x})^{-1}$, respectively $L(1, \ad(r'_{x}))^{-1}$, for all $x\in \X'$.  By purity, $L(1,r, \ad)^{-1}$ does not vanish on $\X^{\cl}$, hence it is invertible in an open neighbourhood $\X''$ of $\X^{\cl}$ in $\X'''$.
\end{proof}


 Let $\X$ be an integral scheme,  $\cF^{\ts}$ be a $\cK_{\X}^{\times}$-module, then we define $\cF^{\ts,-1}$ to be the $\cK_{\X}$-module such that for each open $\cU\subset \X$, 
\beqq
\cF^{\ts,-1}(\cU):= \{ f^{-1}\, |\ \ f\in \cF^{\ts}(\cU)\}\eeqq
  with $\cK_{\X}^{\times}$-action given by $a\cdot f^{-1}=(a^{-1}f)^{-1}$.
 
 

\begin{prop}
\source{universal-p-adic-gross-zagier-formula}
\notready \label{interp QcircS}
\source{universal-p-adic-gross-zagier-formula}
Consider the situation of Lemma \ref{interp Lv}. Let 
$$\pi=\pi(r')$$
 be the $\OO_{\X}[\GL_{2}(F)]$-module attached to $r'$  by   the local Langlands   correspondence in  families of \cite{LLC}, let $\omega\colon F^{\times}\to \OO(\X)^{\times}$ be its central character, and  let $\chi\colon E^{\times}\to \OO(\X)^{\times}$ be a  character  such that $\omega\cdot \chi|_{F^{\times}}=\one$. 
Let $\pi^{\vee}:=\pi(\rho^{*}(1))$ and let $\Pi=\pi \otimes \chi$, $\Pi^{\vee}=\pi^{\vee}\otimes \chi^{-1}$. 
Let $(\Pi\ot_{\OO_{\X}^{\ts}}\Pi^{\vee})^{\ts}$ be the $\OO_{\X}^{\ts}$-submodule of those $f_{3}\ot f_{4}$ such that $(f_{3}, f_{4})\neq 0$. 

 Then there exist: an open subset $\X'\subset \X$ containing $\X^{\cl}$; letting  $\OO:=\OO_{\X'}$, $\cK:=\cK_{\X'} $, an $\OO^{\ts}$-submodule
  $(\Pi\ot_{\OO_{\X}^{\ts}}\Pi^{\vee})^{\ts}$ specialising at all $z\in \X^{cl}$ to the space of $f_{3,z}\ot f_{4,z}$ such that $(f_{3, z}, f_{4, z})_{z}\neq 0$;  and  a map of $\OO$-modules
$$\cQ_{dt}^{} \colon 
(\Pi\otimes_{\OO}\Pi^{\vee})\otimes_{\OO^{\times}}  ({\Pi} \otimes_{\OO^{\times}} {\Pi}^{\vee})^{\ts, -1}
 \to \mathscr{K},$$
satisfying the following properties.
\begin{enumerate}\item
For all $t_{1}, t_{2}\in E^{\times}/F^{\ts}$, $g\in (\GL_{2}(F)\times E^{\times})/F^{}$,
$$\cQ_{dt}\left({ \Pi(t_{1})f_{1}\otimes \Pi^{\vee}(t_{2})f_{2})\over \Pi( g)f_{3}\otimes f_{4}}\right)=
\cQ_{dt}\left( 
{f_{1}\otimes f_{2}\over  f_{3}\otimes \Pi^{\vee}(g^{-1})f_{4}}\right);$$
\item For all $x\in \X^{\cl}$,  
$$\cQ^{}_{dt | x} =Q_{dt}, $$ where $Q_{dt}$ is  the paring on $\Pi_{x}\ot \Pi_{x}^{\vee}$ of   \eqref{Qv sec4}.
\end{enumerate}
\end{prop}
\begin{proof} 
For each $x\in \X^{\rm cl}$, $\pi_{x}$ corresponds to a pure Weil--Deligne representation under local Langlands, hence it is essentially unitarisable (and in fact tempered, see \cite[Lemma 1.4 (3)]{TY}).
Then by \cite[Lemma 5.2.5]{LLC} 
there is an open neighbourhood $\X'$ 
of $\X^{\cl}$ in $\X$ and an invariant  pairing over $\X'$ 
\beq
\label{inv pair family}
\source{universal-p-adic-gross-zagier-formula}
(\ , \, )\colon \pi\otimes \pi^{\vee}\to \OO_{\X'}
\eeq
specialising to   the pairing $(\ , \ )_{x}$ defined by  \eqref{kir pair} at all $x\in \X^{\cl}$.  It induces an invariant pairing $ \Pi\otimes \Pi^{\vee}\to \OO_{\X'}$ still denoted by $(\ , \, )$. 



By Lemma \ref{interp Lv}, up to possibly shrinking $\X'$, we have regular functions on $\X'$ interpolating   $z\mapsto L(1/2, \pi_{z,E}\otimes\chi_{z})^{-1}=L(0, r_{z}|_{W_{E}'} \otimes \chi_{z})^{-1}$ and $x\mapsto L(1, \pi_{x}, \ad)=L(1, r_{x}', \ad)$.

If $E/F$ is split,   \cite[Proposition 5.2.4]{LLC} 
 applied to \eqref{Qv reg} gives an element $\cQ_{(, ),dt}\colon \Pi_{E^{\times}}\otimes \Pi_{E^{\times}}^{\vee}\to \OO_{\X' }$ interpolating $Q_{(, )_{x}}$ for $x\in \X^{\cl}$, and we define
\beq\label{qcircdef}
\source{universal-p-adic-gross-zagier-formula}
\cQ_{dt}\left({f_{1}\ot f_{2}\over f_{3}\ot f_{4} }\right) := {\cQ_{( , ), dt}(f_{1}, f_{2})\over ( f_{3}, f_{4})},
\eeq

If $E/F$ is nonsplit, by the previous discussion we can interpolate all terms occurring in the definition  \eqref{Qpair} (note that the integral there is just a finite sum), to obtain a pairing  $\cQ_{(,)}$ over $\X'$ interpolating $\cQ_{(,)_{x}}$ for $x\in \X^{\cl}$. Then we again define $\cQ^{}$ by \eqref{qcircdef}.
\end{proof}

\subsubsection{Product of local pairings} We  consider the global situation, 
 resuming with the setup of \S\S~\ref{3.2}-\ref{sec: duality}. 
 
 
Let $\Pi:=\Pi^{K^{p}{}',  \ord}_{{{H_{\Sg}'}}}$ over $\X^{(5)}$.
Recall that we have a decomposition 
\beq\label{deco}
\source{universal-p-adic-gross-zagier-formula}
\Pi
\cong (\pi_{\G,\Sg'}(\cV_{\G})\otimes \chi_{\H, {\rm univ},{\Sg'}})
\otimes_{\OO_{\X^{(5)}}}
  \Pi^{K^{p}{}',S,  \ord}_{{H_{\Sg}'}} 
  \eeq
from Theorem \ref{LGC}. 



 Let 
 $
 (\Pi\otimes_{\cK^{\times}_{\X^{(5)}}}\Pi^{\iota})^{\ts}
 \subset  \Pi\otimes_{\cK^{\times}_{\X^{(5)}}}\Pi^{\iota}$ be the 
  $\OO_{\X^{(5)}}^{\times}$-submodule of sections $f_{3}\otimes_{\OO_{\X^{(5)}}^{\times}} f_{4}$ such that $f_{3}\otimes f_{4}\neq 0$ and $(f_{3,v},f_{4,v})_{v}\neq 0$ for each  the pairings $(, )_{v}=$ \eqref{inv pair family}, $v\in \Sg'$. 


 
 
 
 
 \begin{theo} \label{X6} Let $\Pi:=\Pi_{{H_{\Sg}'}}^{K^{p}{}', \ord}$ and $\X^{(5)}$ be as in Proposition \ref{duality}, and let 
\source{universal-p-adic-gross-zagier-formula}
 $
 (\Pi\otimes_{\cK_{\X^{(5)}}}\Pi^{\iota})^{\ts}
 \subset  \Pi\otimes_{\cK_{\X^{(5)}}}\Pi^{\iota}$ be the submodule defined above. 
 Then there exist an open subset $\X^{(6)}\subset \X^{(5)}$ containing $\X^{\cl}$ and, letting  $\OO=\OO_{\X^{(6)}}$, $\cK:=\cK_{\X^{(6)}}$, a map of $\OO^{\ts}$-modules
$$\cQ^{} \colon 
(\Pi\otimes_{\OO}\Pi^{\iota})\otimes_{\OO^{\times}}  ({\Pi} \otimes_{\OO^{\times}} {\Pi}^{\iota})^{\ts, -1}
 \to \mathscr{K}_{\X}$$
satisfying:
\begin{enumerate}\item
For any $t_{1}, t_{2}\in E_{\Sg'}^{\times}/F_{\Sg'}^{\ts}\subset (\GL_{2}(F_{\Sg'})\times E_{\Sg'}^{\times})/F_{\Sg'}^{\ts}$, any $h\in \cH_{S, K_{\Sg'}}$  
and any section 
$$(f_{1}\otimes f_{2}) \otimes  (f_{3}\otimes f_{4})^{-1} \quad 
\text{ of }
\quad
( \Pi\otimes_{\cK} \Pi^{\iota})\otimes_{\cK^{\times}}  ({\Pi} \otimes_{\cK^{\times}} {\Pi}^{\iota})^{\ts, -1}, $$
 we have
$$\cQ^{}\left({\Pi_{\Sg'}(t_{1})f_{1} \otimes \Pi_{\Sg'}^{\iota}(t_{2})f_{2} \over \Pi( h)f_{3} \otimes f_{4}}\right)=
  \cQ^{}\left({ f_{1}\otimes f_{2} \over f_{3} \otimes \Pi^{\iota}(h)f_{4}}\right);$$
 in the left-hand side, $\Pi_{\Sg'}$, respectively $\Pi_{\Sg'}^{\iota}$ denote the actions of the Hecke algebras at $S$ on  $\Pi$, respectively $ \Pi^{\iota}$.
\item For all $x\in \X^{\cl}$,  $$\cQ^{}_{dt |x}=Q^{\ord}, $$  where $Q_{}^{\ord}$  is the restriction of the pairing on $\Pi_{x}^{\vee, \ord}\ot \Pi_{x}^{\vee,\ord}$ of Definition \ref{Q ord def}.
\end{enumerate}
\end{theo}

\begin{proof} By \eqref{Q orddd loc}, \eqref{Q orddd}, \eqref{dec Qord}, and \eqref{deco}, we need to interpolate:
\begin{itemize}
\item the terms $\calL(V_{(\pi, \chi), v}$ for $v\in \Sg$: this is Lemma \ref{interp Lv};
\item the characters $\mu_{v}^{+}$ for $v\vert p$:  this follows form the existence of the filtration  \eqref{decatv} over  an open subset of $\X$.
\item the term $Q_{dt, \Sg'}:=\prod_{v\in \Sg'}Q_{dt, v}$,
 According to the proof of \cite[Theorem 4.4.1]{LLC}, 
the representation $\pi_{\G, {\Sg'}}(\cV_{\G})$ is the maximal torsion-free quotient of $\otimes_{v\in {\Sg'}}\pi_{\G,v}(\cV_{\G})$.  For  sections $f_{i,{\Sg'}}$ that are  images of $\otimes_{v\in {\Sg'}}f_{i,v}$, with $f_{i,v}$ sections of $ \pi_{\G,v}(\cV_{\G})\otimes\chi_{\H, {\rm univ},v}$ if $i=1,3$, or of  $ \pi_{\G,v}(\cV_{\G}^{\iota})\otimes\chi^{-1}_{\H, {\rm univ},v}$ if $i=2,4$, let
$$\cQ^{}_{{\Sg'}}\left({ f_{1, {\Sg'}}\otimes f_{2, {\Sg'}}\over f_{3, {\Sg'}}\otimes f_{4, {\Sg'}}}\right):=
\prod_{v\in {\Sg'}} \cQ_{v}\left( {f_{1,v}\otimes f_{2,v}\over f_{3,v}\otimes  f_{4,v}}\right),$$
where the factors in the right-hand side are provided by Proposition \ref{interp QcircS}. This is
 well-defined  independently of the choices of $f_{i,v}$ as $\mathscr{K}$ is torsion-free.   
\end{itemize}



 
This completes the interpolation of \eqref{Q orddd} into a function $\cQ$, that satisfies properties 1 and 2 by  construction and the corresponding properties from Proposition \ref{interp QcircS}.
\end{proof}
